\chapter{Results}
\label{chap:results}

This chapter evaluates the three reference distributions while holding the MI
difference and standard error fixed. Rejection calibration is the primary
outcome. Coverage, power, validity, and runtime describe the associated
trade-offs.

\section{Rejection calibration across regimes}

Figure~\ref{fig:rejection-calibration} plots actual rejection rate against the
nominal significance level. The diagonal represents perfect calibration. A
curve above the diagonal is liberal and rejects too often; a curve below it is
conservative. Each method's line is the equal-weight mean across the 12 fixed
population pairs in that regime. Shading spans the 10th to 90th percentiles of
the scenario-specific rates.

\begin{figure}[htbp]
  \centering
  \includegraphics[width=\textwidth]{rejection_calibration.png}
  \caption{Rejection calibration under the equal-MI null. Each line is the
  mean across fixed population pairs in one regime, and each band spans the
  10th to 90th percentiles across those pairs. The vertical guide marks
  \(\alpha=0.05\).}
  \label{fig:rejection-calibration}
\end{figure}

The three curves are close in the well-sampled control. The expanded reference
moves the moderate and sparse curves toward the diagonal. The largest visible
change occurs in the ultra-skewed and widespread-sparsity regimes, where the
normal Wald and simple Welch curves are substantially liberal.

Table~\ref{tab:regime-calibration} separates mean FPR from mean absolute
scenario error at \(\alpha=0.05\). Mean FPR describes direction, while mean
absolute error measures scenario-level accuracy without cancellation.

\begin{table}[htbp]
\centering
\caption{Calibration at \(\alpha=0.05\). Error change is the relative
reduction in mean absolute error from normal Wald to expanded Welch; a
negative value denotes deterioration.}
\label{tab:regime-calibration}
\small
\resizebox{\textwidth}{!}{%
\begin{tabular}{@{}lcccccc@{}}
\toprule
Regime & Wald FPR & Simple FPR & Expanded FPR & Wald MAE & Expanded MAE & Change \\
\midrule
Well sampled & 0.05174 & 0.05112 & 0.04913 & 0.00317 & 0.00332 & \(-4.5\%\) \\
Moderate & 0.06158 & 0.06025 & 0.05623 & 0.01158 & 0.00720 & \(37.8\%\) \\
Highly skewed and sparse & 0.05547 & 0.05445 & 0.05198 & 0.00645 & 0.00456 & \(29.3\%\) \\
Ultra-skewed and sparse & 0.07504 & 0.07208 & 0.06305 & 0.02504 & 0.01340 & \(46.5\%\) \\
Widespread sparsity & 0.06621 & 0.06400 & 0.05687 & 0.02087 & 0.01277 & \(38.8\%\) \\
\bottomrule
\end{tabular}%
}
\end{table}

Normal Wald was already calibrated in the well-sampled control. Expanded
Welch changed mean FPR from 0.05174 to 0.04913, but its mean absolute error was
4.5\% larger because some scenarios became mildly conservative. The result is
consistent with the intended asymptotic behaviour: when the component degrees
of freedom are large, little correction is needed.

The moderate regime had normal Wald mean FPR 0.06158. Expanded Welch reduced
it to 0.05623 and reduced mean absolute error by 37.8\%. In the highly skewed
and sparse regime, where both minimum expected counts lay between 1 and 5,
expanded Welch reduced mean absolute error by 29.3\%.

The largest relative improvement occurred in the ultra-skewed regime. Expanded
Welch reduced mean FPR from 0.07504 to 0.06305 and mean absolute error by
46.5\%. The method improved the result but did not restore exact calibration.
The widespread-sparsity boundary showed a similar pattern: mean FPR fell from
0.06621 to 0.05687, while mean absolute error fell by 38.8\%.

At \(\alpha=0.01\), the expanded method reduced mean absolute error by 54.2\%
in the moderate regime, 51.6\% in the highly skewed and sparse regime, 66.3\%
in the ultra-skewed regime, and 48.5\% under widespread sparsity. The heavier
Student tails had a larger relative effect in the more extreme tail.

\section{Variation within regimes}

Regime averages can hide unequal performance between design variants.
Table~\ref{tab:variant-fpr} reports mean FPR across the six table shapes for
each fully specified variant.

\begin{table}[htbp]
\centering
\caption{Mean FPR at \(\alpha=0.05\) by regime variant.}
\label{tab:variant-fpr}
\small
\begin{tabularx}{\textwidth}{@{}Xccccc@{}}
\toprule
Variant & MI & Ratio & Wald & Simple & Expanded \\
\midrule
Well sampled 1 & 0.10 & \(1:1\) & 0.04942 & 0.04882 & 0.04612 \\
Well sampled 2 & 0.15 & \(1:1\) & 0.05407 & 0.05342 & 0.05215 \\
Moderate 1 & 0.10 & \(1:1\) & 0.05772 & 0.05647 & 0.05222 \\
Moderate 2 & 0.15 & \(2:1\) & 0.06545 & 0.06403 & 0.06025 \\
Highly sparse 1 & 0.10 & \(1:1\) & 0.05185 & 0.05077 & 0.04797 \\
Highly sparse 2 & 0.15 & \(2:1\) & 0.05908 & 0.05813 & 0.05598 \\
Ultra-sparse 1 & 0.10 & \(1:1\) & 0.05775 & 0.05638 & 0.05267 \\
Ultra-sparse 2 & 0.15 & \(10:1\) & 0.09233 & 0.08778 & 0.07343 \\
Widespread sparsity 1 & 0.10 & \(1:1\) & 0.06190 & 0.05926 & 0.05057 \\
Widespread sparsity 2 & 0.15 & \(2:1\) & 0.07053 & 0.06874 & 0.06317 \\
\bottomrule
\end{tabularx}
\end{table}

Expanded Welch improved both variants of every difficult regime. The most
severe remaining error was the \(10:1\) ultra-sparse variant. Its mean FPR fell
from 0.09233 to 0.07343 but remained materially above 0.05. This scenario
combines a very small first sample, rare cells, and extreme sample imbalance.
The reference correction cannot remove all higher-order bias and support
instability in that setting.

Figure~\ref{fig:scenario-calibration-comparison} removes the averaging over
shapes and compares absolute FPR error for every population pair. Points below
the diagonal favour the comparator over normal Wald. Simple Welch improved 53
of the 60 scenarios, but most movements were small. Expanded Welch improved 48
scenarios and worsened 12; its improvements were larger in the difficult
regimes, while six of the 12 well-sampled controls became slightly less
accurate. This explains why expanded Welch has the best aggregate error even
though it does not win the largest number of individual scenarios.

\begin{figure}[htbp]
  \centering
  \includegraphics[width=\textwidth]{scenario_calibration_comparison.pdf}
  \caption{Scenario-level absolute FPR error at \(\alpha=0.05\). Each point is
  one fixed equal-MI population pair. Points below the diagonal have lower
  error than normal Wald; colour identifies the sampling regime.}
  \label{fig:scenario-calibration-comparison}
\end{figure}

\section{Overall calibration, coverage, and validity}

Table~\ref{tab:overall-results} aggregates all 60 fixed population pairs.
Expanded Welch had the lowest mean absolute calibration error at each reported
significance level and the coverage closest to 0.95.

\begin{table}[htbp]
\centering
\caption{Overall results across the 60 equal-MI population pairs.}
\label{tab:overall-results}
\begin{tabular}{@{}lccccc@{}}
\toprule
Method & MAE 0.10 & MAE 0.05 & MAE 0.01 & Valid rate & Coverage \\
\midrule
Normal Wald & 0.01716 & 0.01342 & 0.00608 & 0.99489 & 0.93799 \\
Simple Welch & 0.01560 & 0.01194 & 0.00514 & 0.99489 & 0.93962 \\
Expanded Welch & 0.01179 & 0.00825 & 0.00278 & 0.99489 & 0.94455 \\
\bottomrule
\end{tabular}
\end{table}

All methods used the same base validity condition, which required positive
finite first-order variance estimates. In the first four regimes the valid
rate was effectively one. Under widespread sparsity it fell to 0.97447 because
one sampled table could have zero estimated MI influence variance after severe
support loss. Conditioning every method on the same valid table pairs keeps
their reference comparison aligned, but the reduced valid rate is itself a
warning against relying on first-order inference in this regime.

Coverage followed the calibration result. Normal Wald's overall 95\% coverage
was 0.93799, simple Welch increased it to 0.93962, and expanded Welch increased
it to 0.94455. Expanded Welch remained slightly below nominal overall, which
is consistent with the residual liberal behaviour in the most difficult
variants.

\section{Effective degrees of freedom}

Figure~\ref{fig:effective-df-comparison} compares the median total degrees of
freedom assigned by the two Student references in each scenario. All 60 points
fall below the diagonal. The median simple-Welch value was 380.50, whereas the
median expanded value was 71.56; the median ratio was 0.219. The expanded
calculation therefore identifies substantially more uncertainty in the MI
variance estimator than the generic \(n-1\) assignment.

Smaller degrees of freedom do not by themselves establish better inference.
They describe the mechanism by which expanded Welch changes the reference
tails. In the well-sampled scenarios, the median expanded value was still
117.03, for which the Student distribution is already close to normal. In the
ultra-skewed scenarios it fell to 30.27, producing a more material correction.

\begin{figure}[htbp]
  \centering
  \includegraphics[width=0.78\textwidth]{effective_df_comparison.pdf}
  \caption{Median effective degrees of freedom for simple and expanded Welch
  in each fixed null scenario. Both axes are logarithmic; the dashed line
  denotes equality.}
  \label{fig:effective-df-comparison}
\end{figure}

\section{Power}

Table~\ref{tab:power-results} reports power at \(\alpha=0.05\). Power increased
with both the true MI difference and sample size. Simple Welch was nearly
indistinguishable from normal Wald. Expanded Welch had lower power because its
heavier reference tails require a larger \(|T|\) for rejection.

Figure~\ref{fig:power-comparison} shows the same result graphically. Increasing
the sample size from 150 to 600 at fixed \(|\Delta|=0.05\) raises power for all
methods. The ordering between methods remains much smaller than the effect of
sample size or the true MI difference.

\begin{figure}[htbp]
  \centering
  \includegraphics[width=0.92\textwidth]{power_comparison.pdf}
  \caption{Power at \(\alpha=0.05\) for the five alternative scenarios.
  Here \(n=n_P=n_Q\) and \(|\Delta|=|I(P)-I(Q)|\).}
  \label{fig:power-comparison}
\end{figure}

\begin{table}[htbp]
\centering
\caption{Power at \(\alpha=0.05\) for the five alternative scenarios.}
\label{tab:power-results}
\begin{tabular}{@{}cccccc@{}}
\toprule
\(I(P)\) & \(I(Q)\) & \((n_P,n_Q)\) & Wald & Simple & Expanded \\
\midrule
0.05 & 0.07 & \((300,300)\) & 0.0768 & 0.0759 & 0.0689 \\
0.05 & 0.10 & \((150,150)\) & 0.1523 & 0.1498 & 0.1402 \\
0.05 & 0.10 & \((300,300)\) & 0.2775 & 0.2761 & 0.2652 \\
0.05 & 0.10 & \((600,600)\) & 0.5161 & 0.5151 & 0.5063 \\
0.05 & 0.15 & \((300,300)\) & 0.7449 & 0.7437 & 0.7362 \\
\bottomrule
\end{tabular}
\end{table}

Relative to normal Wald, expanded Welch lost 1.02 percentage points of power
on average and at most 1.23 percentage points in these scenarios. The power
cost is small in absolute terms but is not zero. It is the expected trade-off
for correcting a liberal reference with heavier tails.

\section{Computational cost}

Table~\ref{tab:runtime-results} reports the median time for one complete method
call after the count tables had been constructed. Expanded Welch required
approximately 1.72--1.74 times the normal Wald runtime, but all measured times
remained below 0.15 milliseconds.

\begin{table}[htbp]
\centering
\caption{Median end-to-end method runtime over 200 repetitions.}
\label{tab:runtime-results}
\begin{tabular}{@{}ccccc@{}}
\toprule
Shape & \((n_P,n_Q)\) & Wald (ms) & Simple (ms) & Expanded (ms) \\
\midrule
\(2\times2\) & \((200,200)\) & 0.0842 & 0.0991 & 0.1451 \\
\(3\times3\) & \((200,200)\) & 0.0844 & 0.0995 & 0.1456 \\
\(5\times5\) & \((375,375)\) & 0.0837 & 0.0987 & 0.1452 \\
\(8\times8\) & \((960,960)\) & 0.0858 & 0.1009 & 0.1481 \\
\bottomrule
\end{tabular}
\end{table}

The additional arithmetic is negligible in absolute terms. All methods are
deterministic and have \(O(rc)\) complexity. Runtime therefore differs in
constant factors rather than asymptotic order.

\section{Scaled chi-squared mechanism audit}

The expanded derivation uses a scaled chi-squared model for the positive MI
variance estimate. A separate experiment evaluated that working model on 128
population components from 64 fixed scenarios. For each component, 10,000
tables estimated finite-sample moments and an independent 10,000 tables
evaluated distributional fit.

When every candidate distribution received the empirical mean and variance,
the oracle scaled chi-squared model had mean Kolmogorov--Smirnov distance
0.0199, compared with 0.0196 for an oracle normal model and 0.0291 for an
oracle lognormal model. Scaled chi-squared had lower upper-tail error than the
normal model: 0.0045 versus 0.0054 at the 0.05 tail and 0.0026 versus 0.0038 at
the 0.01 tail. It was therefore a plausible tail model but not uniformly the
best whole-distribution family.

The direct population first-order model had mean KS distance 0.1622, showing
that finite-sample error in the predicted moments is more important than the
choice between chi-squared and normal shape families. The implemented plug-in
component degrees of freedom were nevertheless close to empirical
moment-matched values: their median ratio was 0.970. Appendix
\ref{app:additional-validation} reports the full diagnostic.

\begin{figure}[htbp]
  \centering
  \includegraphics[width=\textwidth]{distribution_fit_summary.png}
  \caption{Audit of the scaled chi-squared working model. The left panel
  compares whole-distribution KS distance; the right panel compares absolute
  upper-tail probability error. Oracle models use empirical moments, whereas
  the population first-order model uses the derived moments.}
  \label{fig:scaled-chi-audit}
\end{figure}

\section{Results summary}

The empirical evidence gives a consistent but bounded result. Normal Wald was
already accurate in well-sampled tables. Simple Welch made only a small
change. Expanded Welch improved average calibration, coverage, and tail
behaviour in the moderate and sparse regimes, with the largest gains under
rare cells and unequal sample sizes. It incurred about one percentage point of
average power loss and less than 0.07 milliseconds of additional median
runtime. The most extreme sparse and imbalanced scenario remained liberal,
and widespread support loss produced invalid first-order variance estimates in
approximately 2.6\% of table pairs.
